\documentclass[12pt]{jlreq}
\usepackage{amsmath, amssymb}

\newcommand{\C}{\mathbb{C}}
\newcommand{\R}{\mathbb{R}}
\newcommand{\Z}{\mathbb{Z}}
\newcommand{\N}{\mathbb{N}}
\newcommand{\F}{\mathbb{F}}
\newcommand{\Prob}{\mathbb{P}}
\newcommand{\Exp}{\mathbb{E}}
\newcommand{\dd}{\,d}
\newcommand{\Ker}{\operatorname{Ker}}
\newcommand{\Impart}{\operatorname{Im}}
\newcommand{\Hom}{\operatorname{Hom}}

\begin{document}

無限個の正準変数 \(x = (x_1, x_2, \dots)\)，\(p = (p_1, p_2, \dots)\) の関数 \(f(x,p)\)，\(g(x,p)\) に対し，そのポアソン括弧を，
\[
  \{f,g\} = \sum_{n=1}^\infty \left( \frac{\partial f}{\partial p_n} \frac{\partial g}{\partial x_n} - \frac{\partial f}{\partial x_n} \frac{\partial g}{\partial p_n} \right)
\]
で定義する．\(z\) を不定元とし，母関数 \(\tau_+(z)\)，\(\tau_-(z)\)，\(\eta(z)\) を，
\[
  \tau_+(z) = \exp\left( - \sum_{n=1}^\infty x_n z^n \right),\quad \tau_-(z) = \exp\left( - \sum_{n=1}^\infty p_n z^{-n} \right),
\]
\[
  \eta(z) = \sum_{n\in\mathbb{Z}} \eta_n z^{-n} = \exp\left( \sum_{n=1}^\infty x_n z^n \right) \exp\left( \sum_{n=1}^\infty p_n z^{-n} \right) = \frac{1}{\tau_+(z)\tau_-(z)}
\]
と定める．以下の問に答えよ．

(1) \(\{\eta(z), \eta(w)\}\) を計算することによって，\(m > n\) のとき，次式が成立することを示せ．
  \[
    \{\eta_m, \eta_n\} = \sum_{l=1}^{m-n} \eta_{m-l}\eta_{n+l}.
  \]

以下，\(H = \eta_0\)，\(J = \eta_{-1}\eta_1\)，\(K = \eta_{-1}\eta_{-1}\eta_2 + \eta_{-2}\eta_1\eta_1 - \eta_{-1}\eta_0\eta_2\) とする．

(2) \(\{H,J\} = 0\)，および，\(\{H,K\} = 0\) を示せ．

(3) \(H\) をハミルトニアンとする時間発展の方程式
  \[
    \frac{df}{dt} = \{H,f\}
  \]
  を考える．いま，母関数 \(\tau_+(z)\)，\(\tau_-(z)\) の各係数が (II) にしたがって時間発展するとき，母関数のレベルで
  \[
    \frac{d\tau_-(z)}{dt}\tau_+(z) - \tau_-(z)\frac{d\tau_+(z)}{dt} + H\tau_-(z)\tau_+(z) = 1
  \]
  が成り立つことを，\(\{\eta(w), \tau_+(z)\}\)，\(\{\eta(w), \tau_-(z)\}\) の \(w\) に関する級数展開を用いて示せ．

(4) 常微分方程式 (III) の特殊解を以下のように構成しよう．\(m\) を正の整数とし，\(2m+1\) 個のパラメータ
  \[
    \alpha_1,\, \dots,\, \alpha_m,\, \beta_1,\, \dots,\, \beta_m,\, \omega
  \]
  を用意する．これらを用いて，\(A_k(t), B_k(t), C_k(t), D_k(t)\) (\(0 \le k \le m\)) を，帰納的に，
  \[
    \begin{pmatrix} A_0 & B_0 \\ C_0 & D_0 \end{pmatrix} = \begin{pmatrix} 1 & 0 \\ 0 & 1 \end{pmatrix},\quad
    \begin{pmatrix} A_k & B_k \\ C_k & D_k \end{pmatrix} =
    \begin{pmatrix} A_{k-1} & B_{k-1} \\ C_{k-1} & D_{k-1} \end{pmatrix}
    \begin{pmatrix} 1 & \alpha_k e^{\omega t} z^k \\ \beta_k e^{-\omega t} z^{-k} & 1 \end{pmatrix}
  \]
  と定める．このとき，パラメータ (IV) をうまく選べば，
  \[
    \tau_-(z) = A_m + C_m,\quad \tau_+(z) = B_m + D_m
  \]
  は等高面 \(H = \omega\) に沿って常微分方程式 (III) を満たすという．パラメータ (IV) が満たすべき条件を求めよ．

\end{document}